\documentclass[12pt,a4paper]{article}
\usepackage[brazilian]{babel}
\usepackage[utf8]{inputenc}
\usepackage{amsmath,amssymb,amsthm}
\usepackage{geometry}
\usepackage{enumitem}
\usepackage{hyperref}
\usepackage{xcolor}
\usepackage{graphicx}
\usepackage{fancyhdr}
\usepackage{tcolorbox}
\usepackage{totpages}
\usepackage{titlesec}
\usepackage{etoolbox}
\usepackage{multicol}
\usepackage{multicol}
\setlength{\columnseprule}{0.4pt} % Espessura da linha divisória
\renewcommand{\columnseprulecolor}{\color{borderlight}} % Cor da linha (opcional)

\geometry{top=3.5cm, bottom=2.5cm, left=2.0cm, right=2.0cm, headheight=2.6cm}

\definecolor{primary}{RGB}{10, 45, 90}
\definecolor{secondary}{RGB}{25, 110, 180}
\definecolor{bglight}{RGB}{245, 247, 250}
\definecolor{borderlight}{RGB}{210, 220, 230}
\definecolor{correctgreen}{RGB}{40, 140, 60}

\tcbuselibrary{skins, breakable}

\newtcolorbox{boxfixacao}[1]{
  enhanced,
  colback=bglight,
  colframe=primary,
  arc=3pt,
  boxrule=1pt,
  title={#1},
  fonttitle=\sffamily\bfseries\color{white},
  coltitle=white,
  colbacktitle=primary,
  breakable
}

\newcounter{numquestao}
\setcounter{numquestao}{0}

\newenvironment{questao}[2][]{%
  \refstepcounter{numquestao}%
  \vspace{0.3cm}%
  \begin{tcolorbox}[
    enhanced,
    colback=white,
    colframe=secondary,
    arc=2pt,
    boxrule=0.8pt,
    leftrule=4pt,
    title={\textbf{Questão \thenumquestao} \ifstrempty{#2}{}{--- #2} \hfill \small\textnormal{#1}},
    fonttitle=\sffamily\bfseries\color{white},
    coltitle=white,
    colbacktitle=secondary,
    breakable
  ]%
}{%
  \end{tcolorbox}%
}

\newenvironment{gabarito}{%
  \vspace{0.1cm}%
  \begin{tcolorbox}[
    enhanced,
    colback=bglight,
    colframe=correctgreen,
    arc=2pt,
    boxrule=0.6pt,
    title={\small\sffamily\bfseries Resolução / Gabarito},
    fonttitle=\sffamily\color{white},
    colbacktitle=correctgreen,
    breakable
  ]%
  \small
}{%
  \end{tcolorbox}%
}

\newcommand{\espacoresposta}[1]{%
  \vspace{#1}%
}

\pagestyle{fancy}
\fancyhf{}

\fancyhead[L]{%
  \small\sffamily
  \textbf{IFPE -- Campus Caruaru}
}
\fancyhead[R]{%
  \small\sffamily
  \textbf{Instituto Federal de Pernambuco -- Campus Caruaru} \\
  \textcolor{secondary}{Disciplina:} Cálculo 1 \\
  \textcolor{secondary}{Prof.:} Cleibson José da Silva
}

\fancyfoot[L]{\small\sffamily\color{gray}Lista 4 -- Regras de Derivação}
\fancyfoot[R]{\small\sffamily\color{gray}Página \thepage\ de \ref{TotPages}}

\renewcommand{\headrulewidth}{0.8pt}
\renewcommand{\headrule}{{\color{primary}\hrule width\headwidth height \headrulewidth}}
\renewcommand{\footrulewidth}{0.4pt}

\hypersetup{
    colorlinks=true,
    linkcolor=secondary,
    urlcolor=secondary,
    citecolor=primary
}

\begin{document}

\begin{center}
  {\LARGE\bfseries\color{primary} Lista 4: Regras de Derivação} \\[0.1cm]
  {\small\sffamily\color{gray} Exercícios sobre regras de derivação}
\end{center}

\vspace{0.2cm}


\section*{Lista 04 - Regras de Derivação}

\begin{enumerate}
    \item Calcule as derivadas:
    \begin{enumerate}
        \item $y = (2x^3 - x^2 + 1)(x - 4)$
        \item $y = \dfrac{3x^2 + 1}{x^3 - 2}$
        \item $y = e^{x^2}$
        \item $y = \ln(\sin x)$
    \end{enumerate}

    \item Aplique a regra da cadeia e interprete o resultado:
    \begin{enumerate}
        \item $y = \sin(3x^2 + 2x)$ — qual o significado da derivada em termos da variação de $x$?
        \item $y = \sqrt{1 + e^{2x}}$ — como a função cresce em comparação com $e^x$?
    \end{enumerate}

    \item Uma partícula move-se segundo $s(t) = t e^t$, onde $s$ é a posição em metros e $t$ em segundos.
    \begin{enumerate}
        \item Determine a velocidade instantânea em $t = 1$.
        \item Interprete o valor encontrado em termos físicos.
    \end{enumerate}

    \item O custo de um produto é dado por $C(x) = x^3 - 6x^2 + 9x$. Determine:
    \begin{enumerate}
        \item $C'(x)$, a taxa de variação marginal do custo.
        \item O ponto em que o custo cresce mais rapidamente.
    \end{enumerate}

    \item Uma curva é dada por $y = e^x \cos x$. Determine $y'$ e descreva qualitativamente como o comportamento oscilatório da função se mistura com o crescimento exponencial.

   \item Um corpo se move segundo $s(t) = t^3 - 6t^2 + 9t + 2$ (SI).  
    Determine:
    \begin{enumerate}
        \item As expressões para velocidade e aceleração.  
        \item Os instantes em que o corpo muda o sentido do movimento.  
        \item O intervalo de tempo em que o corpo está se afastando da origem.
    \end{enumerate}

    \item A resistência de um fio condutor depende da temperatura pela expressão $R(T) = R_0(1 + 0{,}004T)$.  
    Calcule $R'(T)$ e interprete o resultado para $R_0 = 10\,\Omega$ em $T = 50^\circ$C.

    \item Uma peça cilíndrica deve ter volume constante $V = 1000\,\text{cm}^3$.  
    Determine as dimensões (raio e altura) que minimizam a área da superfície total do cilindro.  
    (Sugestão: escreva $A(r) = 2\pi r^2 + \frac{2V}{r}$ e derive.)

    \item Uma barra metálica aquecida uniformemente sofre deformação segundo $y(x) = x^3 - 6x^2 + 9x$.  
    Determine os pontos onde a deformação é máxima e mínima e discuta o significado mecânico desses pontos.

    \item O rendimento de uma turbina hidráulica pode ser modelado por $\eta(v) = v e^{-0{,}1v}$, onde $v$ é a velocidade do fluido (m/s).  
    Determine a velocidade que maximiza o rendimento da turbina e o valor máximo de $\eta$.



\end{enumerate}

  
\end{document}